%%%%%
%
\section{FSSP: Rapport de m\'elange en eau liquide, pour 
des classes lin\'eaires en diam\`etre}

\itemstyle{Op\'erateur~:} \operfmt{fssp\_lwc\_lin\_tab}
\label{OP-fssp_lwc_lin_tab}
\index{fssp\_lwc\_lin\_tab}

%%%%%OBJET%%%%
\itemstyle{Objet~:}
Calcul du rapport de m\'elange en eau liquide (g/m3)
� partir d'un tableau de comptages issu d'un granulom�tre
et d'une concentration totale, pour le cas o�
le diam�tre moyen dans chaque classe est
une fonction lin�aire du num�ro de classe.

%%%%%ARGUMENTS%%%%
\bigskip
%%%%%ENTREES%%%%
\begin{tabarg}
\multicolumn{3}{l}{\titretab {Entr\'ees}}\\
\hline
$C$ & Vecteur, \texttt{float[15]} & Tableau des comptages ()\\
$Conc$ & Vecteur, \texttt{float} & Concentration totale de gouttelettes (cm$^{-3}$)\\
$a$ & Coeff., \texttt{float} & Terme de proportionnalit� du calcul du diam�tre moyen (largeur de classe) ($\mu m$)\\
$b$ & Coeff., \texttt{float} & Terme constant du calcul du diam�tre moyen ($\mu m$)\\
\end{tabarg}

\medskip
%%%%%SORTIES%%%%
\begin{tabarg}
\multicolumn{3}{l}{\titretab {Sorties}}\\
\hline
LWC & Vecteur, \texttt{float} & Rapport de m\'elange en eau liquide (gram/m$^{3}$)\\
\end{tabarg}


%%%%%FORMULE%%%%
\itemstyle{Formule}

%%fssp_lwc_lin_tab : calcul du rapport de m\'elange en eau liquide 
% a partir d'une concentration totale et tableau  de comptage
% dans lequel le diam�tre moyen mesur� est fonction lin�aire
% de la classe.
\begin{align*}
\rho_w &= 1000 kg/m^3 \\
\phi_k &= a \cdot k + b  \qquad \text{(diam\`etre m\'edian classe $k$)}\\
LWC &=  \rho_w .\frac{\pi}{6}. Conc .
\frac{\displaystyle\sum_{k=1}^{15} \phi_k^3.C_k} 
     {\displaystyle\sum_{k=1}^{15} C_k} 
\end{align*}


%%%%% COMMENTAIRE %%%
\itemstyle{ Commentaire~: }
Si les comptages sont disponibles sous la forme de
15 vecteurs correspondant chacun � une classe, 
utiliser l'op�rateur 
\hyperref[OP-fssp_lwc_lin]{\operfmt{fssp\_lwc\_lin}}.

%%%%%AUTEUR%%%%
\itemstyle{Auteur~:}
CNRM/GMEI/TRAMM

%%%%%REFERENCE%%%%
\itemstyle{R\'ef\'erences~:}
Bulletin NCAR-RAF Nr 24 \cite{RAF-24}.
